\section{Conclusion}
\label{sec:conclusion}

This study presents a CausalRAG pipeline for Vietnamese criminal-law
question answering, in which the central retrieval object is a graph of
normative legal condition--effect relations and the verification object is
an executable rule model. The pipeline combines rule normalization, causal
memory, multi-hop event/rule retrieval, causal-path ranking, query-aware
verification, and extractive answer generation with citations.

LegalSCM is constructed as a deterministic three-valued structural causal
model of legal norms. Within the scope of the encoded graph, the model can
validate a factual path, perform a hard intervention on a mediator, and
re-infer the outcome. In parallel, path ablation is retained as a
reachability baseline after node deletion and is not treated as an
implementation of the $\doop(\cdot)$ operator.

On the $250$-question benchmark, the pipeline achieves Rule Recall@5 of
$69.93\%$, Event Recall@5 of $74.66\%$, Top-1 Exact Path of $35.20\%$,
Oracle Exact Path of $71.20\%$, Verification Accuracy of $92.00\%$, Answer
Token F1 of $68.69\%$, ROUGE-L F1 of $67.67\%$, and Citation F1 of
$62.86\%$. The $36$-percentage-point gap between Top-1 and Oracle Exact Path
indicates that path ranking is the principal bottleneck. Reverse reasoning
is the weakest group, whereas the linear path metric does not appropriately
capture branch and convergence structures.

The paired ablation records no difference in any quality metric between
structural SCM and path ablation. This is a null quality result, not evidence
of superior structural verification or of theoretical equivalence between
the two methods. Structural SCM provides intervention semantics and more
detailed world-state traces, but requires approximately $3.23$ times more
mean runtime per sample. The current benchmark, particularly its $20$
direct-edge negative samples, is not sufficiently sensitive to quantify any
accuracy benefit of hard intervention.

The study's main contribution therefore lies in constructing and evaluating
a pipeline with clearly bounded claims, providing traceable paired
artifacts, and precisely identifying the remaining bottlenecks. The decisive
next step is not merely to increase an aggregate score, but to construct a
balanced, expert-validated structural-counterfactual benchmark while
improving query-aware path ranking and reverse retrieval. Until these steps
are completed, the system should be regarded as a research prototype for
supporting analysis, not as a legal advisory or decision-making tool.
